\subsection{The rational three-state comparison at order thirteen}

\begin{lemma}[Proved: two column-character covectors on at most five rows]
\label{CP2:A-two-covector-affine-row-origin}
Let $M$ have rectangular minimal band $I\times\{0,1\}$ with
$|I|\leq5$, let $V=\widetilde{\mathbb QI}$, and let
$\chi(m)=\tau_m(1)-\tau_m(0)$ be the old column character.
Put $A=\mathbb QM/\mathbb QK$. If
\[
 \dim\operatorname{Hom}_A(V,\mathbb Q_\chi)\geq2,
\]
then every finite ordinary right-act cone over $M$ has zero
rational degree-two products. No homology bound is required.
\end{lemma}
\begin{proof}
Choose independent covectors $\ell_1,\ell_2$ and a base row
$i_0$. The map
\[
 f(i)=\bigl(\ell_1([i]-[i_0]),\ell_2([i]-[i_0])\bigr)
\]
has a finite image $F$ affinely spanning $\mathbb Q^2$.
Every actor satisfies
$f(\sigma_m i)=\chi(m)f(i)+b_m$ for a vector $b_m$.
If $\chi(m)=1$, its translation preserves the finite set
$F$, so $b_m=0$. If $\chi(m)=-1$, it is a central
reflection of $F$. Every such reflection is the same,
because the composite of two is a translation preserving
$F$ and hence the identity.

If there is no negative actor, $|F|\geq3$ and $|I|\leq5$
ensure a singleton fiber. Its row is fixed by every actor
with nonconstant column map. If a negative actor exists,
central symmetry and affine dimension two give $|F|\geq4$.
There is then at most one nonsingleton fiber. A centrally
symmetric spanning set has at least two noncentral opposite
pairs, so at least one opposite pair has singleton fibers.
The midpoint of their two row basis vectors is fixed by
every nonconstant-column actor. In either case we have a
common rational affine row origin of augmentation one.
Lemma~\ref{CP2:A-act-common-affine-origin} proves the result.
\end{proof}

\begin{proposition}[Proved: rational three-state cones through order thirteen]
\label{CP2:A-thirteen-three-state-rational}
An ordinary right-act cone over a monoid of order at most
thirteen, on at most three states, cannot have the rational
homology and cohomology ring of $\CP^2$.
\end{proposition}
\begin{proof}
Orders through twelve are
Theorem~\ref{CP2:A-bounded-twelve-acts}. Suppose the source
has order thirteen. The relative-product and two-row
ordinary-act obstructions require nonzero $H^2(BM;\mathbb Q)$,
at least three minimal rows and at least two minimal columns.
If the minimal group is nontrivial, the only possibility
is the $3\times2$ Rees ideal with group of order two and
one outside identity. Its actual averaged quotient has
pure $3\times2$ band with identity, projective row module,
and unconditional zero products on every act cone. The
actual rational $H^2$ comparison of
Theorem~\ref{CP2:A-minimal-group-collapse-collage-H2} transfers
that obstruction. No boundedness of its target is used.

The minimal ideal is therefore a band. If some band actor
is nonconstant, the constant-action quotient has at most
two states; its second products vanish and its pullback
surjects on $H^2$ by Lemma~\ref{CP2:A-constant-quotient-H2}.
We may assume constant band action and exactly three
states. Compression is an actual homotopy equivalence
to a monoid with minimal band $I\times S$ and the same
contracted algebra $A$.

Use the complete seven-shape list in
Proposition~\ref{CP2:A-thirteen-five-state-single-shape}.
The $3\times2$ source is excluded unconditionally by
Theorem~\ref{CP2:A-three-row-seven-outside-acts}; the
$5\times2$ and $6\times2$ sources use the bounded
two-column/four-outside theorem. The $3\times4$ and
$4\times3$ sources have contracted algebra $\mathbb Q$
and projective row module, so all second products vanish.
For source shape $3\times3$ with four outside elements,
the compressed row and column modules both have dimension
two over a four-dimensional algebra. If $A$ is local,
Lemma~\ref{A:local-four-small-target} excludes the required
Tor dimensions $(1,0,1,0)$; if nonlocal,
Lemma~\ref{A:nonlocal-four} excludes bounded nonzero
second Tor. This uses rational Tor only, not the actual
monoid minimum.

It remains that the source band is $4\times2$ and
$\dim A=5$. Its row module $V$ has dimension three and
the compressed column module $W$ has dimension two.
If $A$ is local, Corollary~\ref{CP2:A-local-two-column-acts}
kills all products on the original act cone. For nonlocal
$A$, Lemma~\ref{A:nonlocal-five-target} leaves the local
four-dimensional product and the two exceptional algebras.
In $B\times\mathbb Q$, all positive Tor comes from the
$B$ summands, whose tensor product must have dimension
one. Applying Lemma~\ref{A:local-four-small-target} with
the sides exchanged excludes them: their dimensions are
at most two and three. The lemma applies to $B^{\rm op}$,
which is still local of dimension four.

For the line algebra, the only pair of dimensions
$(\dim V,\dim W)=(3,2)$ with Tor triple $(1,0,1)$ is
\[
 V=S_3\oplus Ae_2,\qquad W=T_1\oplus T_3.
\]
Indeed the degree-two condition forces one $S_3$ and one
$T_1$, and degree-one vanishing excludes $S_2,T_2$.
The right dimension then forces $W=T_1\oplus T_3$;
among $2S_1\oplus S_3$, $S_3\oplus Ae_2$ and
$S_3\oplus Ae_3$, only the displayed module gives
tensor dimension one. Its balanced form is supported
on the single $T_3\times S_3$ coordinate and has matrix
rank one. Lemma~\ref{A:rank-one-square} makes its square
zero on the actual compression monoid, a contradiction.

For the Nakayama algebra, the formulas in
Lemma~\ref{CP2:A-five-dimensional-three-four-target}
for the first two positive Tor dimensions apply without
change. The sole pair of dimensions $(3,2)$ with Tor
triple $(1,0,1)$ is
\[
 V=S_2\oplus Ae_2,\qquad W=R=e_1A/\mathbb Qc.
\]
The only simple quotient character of $V$ is $S_2$,
and its multiplicity in the top is two. The old column
character $\chi$ is a one-dimensional simple character.
The nonzero group $H^2(BM;\mathbb Q)$ is naturally
$\operatorname{Hom}_A(V,\mathbb Q_\chi)$, by the
band--Tor balanced description. Thus $\chi=S_2$ and
this Hom space has dimension two. The preceding affine
origin lemma applies to the four original rows and kills
all products on the original act cone.

The exact program
\path{monoid/code/five_dimensional_three_two_pairs.py}
reconstructs the literal module actions and
resolution ranks used in this dimension check. It tests
all $128$ line pairs and all $45$ Nakayama pairs, obtaining
exactly the two displayed survivors. The four-rule
derivation from the accepted pair checker, all tables
and completed output are under
\path{two_object19/ordinary13_five_algebra/}.
The finite calculation only confirms the complete
indecomposable arithmetic just given; it does not assert
a census of thirteen-element source monoids.
\end{proof}

\begin{corollary}[Proved: the remaining order-two opposite budget]
\label{CP2:A-oriented-nineteen-thirteen-two-four}
No two-object category in the complete oriented budget
$(13,2,4)$ has type $\CP^2$.
\end{corollary}
\begin{proof}
There are exactly two monoids of order two. For
$\{1,e\}$ with $e^2=e$, one has $eN=eNe=\{e\}$.
The sink form of
Theorem~\ref{CP2:A-extreme-idempotent-corner-reduction}
gives an actual homotopy equivalence to a category with
at most eighteen arrows, contradicting the adopted
eighteen-arrow exclusion.

For $N=C_2$, a nontrivial action on the four connecting
states has at most three orbits. Exact rational group
coinvariants identify the actual orbit-quotient functor
with a rational homology and cohomology-ring equivalence
to an ordinary $M$-act cone on those orbits. The preceding
proposition excludes the rational projective-plane pattern.

If the $C_2$ action is trivial, the ordinary act cone
on the original four states is an actual categorical
retract, and its collapse is a rational homology
equivalence. Every integral homology group of the retract
is therefore the corresponding direct summand of that
of a projective plane, with ranks one in degrees two
and four and zero in every other positive degree.
The induced isomorphism on integral degree two and the
actual cup pullback give a primitive square. Its
fundamental group retracts the trivial group. Projective-plane
recognition makes this an actual $\CP^2$ model, but it
has only eighteen arrows. This last contradiction uses
an actual retract; no rational equivalence is promoted
to an integral homotopy equivalence by itself.
\end{proof}
